\section{Conclusion}
\label{sec:conclusion}

This paper studied passenger-response-aware bidirectional meeting-point DRT service
design for many-to-many operations. The framework links candidate pickup and dropoff
meeting-point generation, single-offer binary-logit passenger response, and
rolling-horizon dispatch. It is designed to evaluate operational trade-offs: whether
meeting-point consolidation can reduce routing intensity for served trips under tested
synthetic service-design conditions, and what happens to coverage, choice rejection,
feasibility rejection, and simulated passenger-type heterogeneity when that
consolidation is introduced.

The contribution is conditional. Under the tested simulation settings,
bidirectional meeting-point consolidation can produce lower routing intensity among
served trips, but this signal must be interpreted with the accompanying coverage and
passenger-type trade-offs. The full passenger-response-aware design is therefore not a
universal substitute for door-to-door service. It is an operational design option whose
usefulness depends on demand density, fleet deployment, walking tolerance, and the
quality of offers presented to passengers.

The paper also clarifies evidence boundaries for future manuscript and package work.
Primary behavioral claims should be based on formal Phase 6 evidence and should report
the relevant denominators. Matched-coverage, fixed-accepted-set, robustness,
passenger-type, Gamma, Beijing-inspired, and algorithm materials are diagnostic,
robustness, monitoring, or limitation evidence unless Phase 4 verifies stronger
source-specific wording. Gamma remains post-hoc welfare or sensitivity accounting. The
Beijing-related evidence remains synthetic. The MILP remains a simplified ex-post
diagnostic over fixed accepted sets.

Several limitations remain. The passenger-response parameters are simulation inputs, not
empirical calibration from real users. The scenarios are synthetic or Beijing-inspired
rather than public-data validations. Passenger types represent service-sensitivity
ranges, not real demographic groups. The routing diagnostic is simplified and fixed-set,
not a complete dynamic exact model. These limitations are important because they define
what the current evidence can and cannot support.

Future work should address these boundaries directly. Real public-data ingestion and
field calibration would support empirical passenger-response estimates. Endogenous
Gamma behavior would require a model in which welfare accounting affects offer or
routing decisions by design. Stronger exact methods would be needed to compare the full
online behavioral process rather than a fixed accepted set. More detailed passenger data
would also allow passenger-type monitoring to be replaced by empirically grounded
equity and burden analysis.
